\documentclass[11pt, a4paper]{article}
\usepackage[a4paper, margin=2.5cm]{geometry}
\usepackage{fontspec}
\usepackage[italian, bidi=basic, provide=*]{babel}
\babelprovide[import, onchar=ids fonts]{italian}
\babelfont{rm}{Noto Sans}
\usepackage{tikz}
\usepackage{amsmath}
\usetikzlibrary{arrows.meta, calc, backgrounds}

\begin{document}

\begin{figure}[htbp]
  \centering
  \begin{tikzpicture}[>=Stealth, font=\small]

    % Stili personalizzati per elementi grafici
    \tikzset{
      gridline/.style={draw=gray!25, line width=0.3pt, dashed},
      badge/.style={draw=gray!40, fill=gray!10, rounded corners=3pt, inner xsep=8pt, inner ysep=4pt, font=\small\sffamily},
      forcearrow/.style={line width=1.3pt, ->},
      posarrow/.style={line width=0.8pt, ->, color=black!80}
    }

    % ==========================================
    % CASO 1: Repulsivo (q q' > 0)
    % ==========================================
    \begin{scope}[shift={(0,0)}]
      % Griglia tecnica di sfondo
      \draw[gridline] (-2.0,-1.1) grid (7.8,1.1);

      % Anelli concentrici (rappresentazione infografica di campo/potenziale)
      \foreach \r in {0.25, 0.45, 0.65} {
        \draw[gray!45, line width=0.4pt, dashed] (0,0) circle (\r);
        \draw[gray!45, line width=0.4pt, dashed] (4.2,0) circle (\r);
      }

      % Vettore Posizione r
      \draw[posarrow] (0,0) -- (4.2,0) node[midway, below=3pt] {$\vec{r}$};

      % Vettori Forza F' e F con prolungamento tratteggiato
      \draw[forcearrow, color=orange!90!red] (0,0) -- (-1.4,0) node[above=2pt] {$\vec{F}'$};
      \draw[dashed, color=orange!90!red, line width=0.8pt] (-1.4,0) -- (-1.8,0);

      \draw[forcearrow, color=orange!90!red] (4.2,0) -- (5.6,0) node[above=2pt] {$\vec{F}$};
      \draw[dashed, color=orange!90!red, line width=0.8pt] (5.6,0) -- (6.0,0);

      % Particelle / Cariche
      \filldraw[fill=orange!80!red, draw=black!80, line width=0.6pt] (0,0) circle (3.2pt);
      \filldraw[fill=purple!70!blue, draw=black!80, line width=0.6pt] (4.2,0) circle (3.2pt);

      % Etichette P, q, P', q'
      \node[above=4pt, font=\small\bfseries] at (0,0) {$P$};
      \node[below=4pt] at (0,0) {$q$};
      \node[above=4pt, font=\small\bfseries] at (4.2,0) {$P'$};
      \node[below=4pt] at (4.2,0) {$q'$};

      % Badge identificativo della condizione
      \node[badge] at (6.8,0) {$q\,q' > 0$};
    \end{scope}

    % ==========================================
    % CASO 2: Attrattivo (q q' < 0)
    % ==========================================
    \begin{scope}[shift={(0,-2.5)}]
      % Griglia tecnica di sfondo
      \draw[gridline] (-2.0,-1.1) grid (7.8,1.1);

      % Anelli concentrici
      \foreach \r in {0.25, 0.45, 0.65} {
        \draw[teal!40, line width=0.4pt, densely dotted] (0,0) circle (\r);
        \draw[teal!40, line width=0.4pt, densely dotted] (4.2,0) circle (\r);
      }

      % Vettore Posizione r
      \draw[posarrow] (0,0) -- (4.2,0) node[midway, below=3pt] {$\vec{r}$};

      % Vettori Forza F' e F (elevati per evitare sovrapposizioni)
      \draw[forcearrow, color=teal!80!black] (0,0.3) -- (1.4,0.3) node[midway, above=1pt] {$\vec{F}'$};
      \draw[dotted, color=teal!60, line width=0.6pt] (0,0) -- (0,0.3);

      \draw[forcearrow, color=teal!80!black] (4.2,0.3) -- (2.8,0.3) node[midway, above=1pt] {$\vec{F}$};
      \draw[dotted, color=teal!60, line width=0.6pt] (4.2,0) -- (4.2,0.3);

      % Particelle / Cariche
      \filldraw[fill=cyan!80!black, draw=black!80, line width=0.6pt] (0,0) circle (3.2pt);
      \filldraw[fill=blue!80!black, draw=black!80, line width=0.6pt] (4.2,0) circle (3.2pt);

      % Etichette P, q, P', q'
      \node[above=4pt, font=\small\bfseries] at (0,0) {$P$};
      \node[below=4pt] at (0,0) {$q$};
      \node[above=4pt, font=\small\bfseries] at (4.2,0) {$P'$};
      \node[below=4pt] at (4.2,0) {$q'$};

      % Badge identificativo della condizione
      \node[badge] at (6.8,0) {$q\,q' < 0$};
    \end{scope}

  \end{tikzpicture}
  \caption{Interazione elettrostatica tra due cariche puntiformi $q$ e $q'$ nei casi repulsivo ($q\,q' > 0$) ed attrattivo ($q\,q' < 0$).}
  \label{fig:legge_coulomb_infografica}
\end{figure}

\end{document}